% Flavor catalog per-process draft.
% process_id: L006
% family: charged_lepton
% owner_agent_id: PKA-L006
% writer_agent_id: null
% checker_agent_id: null
% opus_signoff_id: null
% Canonical metadata sidecar: L006.yaml
% Status is controlled by the YAML sidecar status_history, not this comment.

\section{L006: Muonium--Antimuonium Conversion}

\subsection*{Process}
\textbf{Standard notation:} \(P_{M\bar M}\) and \(G_C/G_F\), where
\(M=(\mu^+ e^-)\) and \(\bar M=(\mu^- e^+)\).  This charged-lepton entry covers
spontaneous muonium-to-antimuonium conversion, a \(\Delta L_\mu=-\Delta L_e=2\)
lepton-flavor-violating oscillation process, not a dipole decay.

\subsection*{PDG / equivalent value(s)}
The sidecar block
\texttt{pdg\_or\_equivalent.pdg\_effective\_coupling\_limit}, from the 2026
PDG muon listing datablock S004MC, gives the canonical effective-coupling
limit
\[
  G_C/G_F < 0.0030 \qquad (90\%\ {\rm C.L.})
\]
for the stated \((V-A)\times(V-A)\) four-lepton effective Lagrangian.  The same
PDG datablock, mirrored in
\texttt{pdg\_or\_equivalent.primary\_current\_probability\_limit}, records the
underlying MACS/PSI probability limit from Willmann et al.,
\[
  P_{M\bar M} < 8.3\times 10^{-11}\qquad (90\%\ {\rm C.L.})
\]
in a \(0.1~{\rm T}\) magnetic field.  The local canonical snapshot is
\texttt{flavor\_catalog/references/L006/pdg2026\_muonium\_antimuonium\_s004mc.txt};
the primary MACS article snapshot is
\texttt{flavor\_catalog/references/L006/willmann1999\_macs\_arxiv\_hepex9807011.txt}.
These are upper limits, not measurements of a nonzero conversion rate.

\subsection*{Relevance to the RS / anarchic-flavor pipeline}
Muonium conversion probes a four-lepton contact structure
\((\bar\mu\Gamma e)(\bar\mu\Gamma e)\) that violates muon and electron family
numbers by two units.  In an RS lepton extension, such operators could arise
from flavor-misaligned heavy neutral or doubly charged states, KK-sector
lepton-current misalignment, or other non-dipole lepton interactions.  It is
therefore complementary to L001: the present repo checks only
\(\mu\to e\gamma\), while L006 would require a separate \(\Delta L=2\)
four-lepton observable.

\subsection*{Post-2008 developments}
Relative to the \texttt{CFW2008} era, the experimental bound itself has not
changed: PDG still uses the 1999 MACS/PSI result.  The main post-2008 movement
is prospective.  The sidecar \texttt{prospects} entries record the Snowmass
2021 statement that MACE is intended to improve probability sensitivity by
more than two orders of magnitude and the 2024 MACE conceptual-design target
beyond the \(10^{-13}\) level.  On the theory side,
\texttt{post\_2008\_theory\_context} records recent EFT work separating
\(\Delta L_\mu=-\Delta L_e=2\) operators from ordinary one-unit LFV and noting
that muonium conversion probes only part of that operator class.

\subsection*{Constraint validity and model dependence}
The experimental signature is clean and lepton-sector-only: a confirmed signal
would be physics beyond the Standard Model.  The quoted \(G_C/G_F\) limit is
not operator-universal, however; it is tied to the PDG \((V-A)\times(V-A)\)
normalization.  Any catalog-to-code translation must choose the Lorentz basis,
magnetic-field treatment, and spin-state normalization before comparing a model
Wilson coefficient to the MACS probability limit.

\subsection*{Code coverage in this repo}
\textbf{NO.}  A targeted grep for
\texttt{muonium|antimuonium|Muonium|Antimuonium|MACE|MACS} over
\texttt{quarkConstraints/}, \texttt{qcd/}, \texttt{flavorConstraints/},
\texttt{neutrinos/}, \texttt{yukawa/}, \texttt{warpConfig/}, \texttt{solvers/},
\texttt{scanParams/}, and \texttt{tests/} returned no hits.  The adjacent
charged-lepton LFV implementation is the \(\mu\to e\gamma\) dipole checker at
\texttt{flavorConstraints/muToEGamma.py}:75, called from
\texttt{scanParams/scan.py}:524; it does not cover muonium conversion.

\subsection*{Implementation difficulty}
\textbf{MEDIUM.}  L006 needs a new \(\Delta L_\mu=-\Delta L_e=2\) four-lepton
operator convention and a bound-state conversion-probability wrapper.  It does
not require lattice inputs or hadronic long-distance calculations for a first
catalog-to-code implementation, but a production version should handle
operator-dependent magnetic-field and spin factors explicitly.

\subsection*{Key references}
Process-local keys before bibliography consolidation:
\texttt{PDG2026\_MuoniumAntimuoniumS004MC}, \texttt{Willmann1999\_MACS},
\texttt{Bai2022\_SnowmassMuonium}, \texttt{Bai2024\_MACECDR},
\texttt{HeeckSokhashvili2024\_DeltaLTwo},
\texttt{AgasheBlechmanPetriello2006\_RSLFV}, and \texttt{CFW2008}.
